\section{The complete conductor of an original arithmetic heat collision}
\label{sec:hc}
\newcounter{hc}\renewcommand{\thehc}{HC\arabic{hc}}
\newcommand{\HC}{\refstepcounter{hc}\par\medskip\noindent\textbf{\thehc.}\ }

\HC\textbf{Original objects and the purpose of the comparison.}
Use precisely the Rodgers--Tao family and arithmetic coordinate of
EH1 and EP1--17. Fix any actual real zero $x_0$ of $H_0$, of order
$m\ge1$. Put $\xi=Z-x_0$, $\mathcal O=\mathbb C\{t\}$ and
$\mathcal K=\operatorname{Frac}\mathcal O$. The full coefficients
$a_n$, $c=a_{m+1}/a_m$ and $d=a_{m+2}/a_m$ remain those of EP1.
The existence of a multiple zero is not asserted. Every calculation
below applies to every actual zero, including $m=1$.

EP2--7, proved from the original heat equation, give all the local
branches after the exact base change $t=u^2$:
\[
 \theta_j(u)=x_j(u^2)-x_0
   =r_j u+2cu^2+r_j(2d-c^2)u^3+O(u^4),\qquad
 \theta_j(-u)=\theta_{m+1-j}(u).
 \tag{HC1}
\]
The equality of branches follows by uniqueness in the analytic
implicit function theorem: the same entire function is evaluated
at $u^2=(-u)^2$, and the leading coefficient changes from $r_j$
to $-r_j=r_{m+1-j}$. Each $\theta_j$ here is the complete analytic
germ, not the displayed truncation.

Let $r=\lfloor m/2\rfloor$ and $\varepsilon=m-2r\in\{0,1\}$.
Index the positive Hermite roots by $b_1<\cdots<b_r$ and let
$\theta_a$ be the corresponding branches. Set
\[
 \begin{split}
 E_a(t)&=\frac{\theta_a(u)+\theta_a(-u)}2,\\
 R_a(t)&=\frac{\theta_a(u)-\theta_a(-u)}{2u},\\
 Q_a(t,\xi)&=(\xi-E_a(t))^2-tR_a(t)^2.
 \end{split}
 \tag{HC2}
\]
Evenness makes $E_a,R_a$ elements of $\mathcal O$;
$R_a(0)=b_a\ne0$ and $E_a(t)=2ct+O(t^2)$.
When $\varepsilon=1$, the central branch is even in $u$;
write it $\theta_0(t)$ and set $Q_0=\xi-\theta_0(t)$.
It has $\theta_0(t)=2ct+O(t^2)$.

\HC\textbf{The full finite algebra in the original coordinate.}
Define
\[
 W(t,\xi)=\prod_{a=1}^r Q_a(t,\xi)\,Q_0(t,\xi)^{\varepsilon},
 \qquad A=\mathcal O[\xi]/(W).
 \tag{HC3}
\]
The convention omits $Q_0$ when $\varepsilon=0$. This monic
polynomial has degree $m$, has all the actual nearby zeros, and
$W(0,\xi)=\xi^m$. Thus $1,\xi,\ldots,\xi^{m-1}$ is an
$\mathcal O$-basis and $A/tA=\mathbb C[\xi]/(\xi^m)$.

For completeness the factorization retains the full analytic unit:
\[
 H_t(x_0+\xi)=U(t,\xi)W(t,\xi),\qquad U(0,0)=a_m\ne0.
 \tag{HC4}
\]
Here is a direct construction. Take a small circle around $\xi=0$
containing no other zero of $H_0$, and shrink the $t$ disc so its
boundary has no zero of $H_t$. EP4's argument-principle calculation
counts exactly $m$ zeros inside. On that boundary $H_t/W$ is
jointly analytic. Its Cauchy integral in the circle defines a
jointly analytic function in the interior. For nonzero sufficiently
small $t$ the roots are simple by HC1, so the removable-singularity
theorem identifies this Cauchy integral with $H_t/W$ everywhere.
The equality extends to $t=0$ by continuity. At $(0,0)$ its value
is $a_m$; on a smaller neighbourhood it is a unit. This proves HC4
without discarding any coefficient of $H$.

\HC\textbf{An exact inclusion into the separated branch algebra.}
For each paired component introduce its own parameter $u_a$ with
$u_a^2=t$. Define
\[
 B=\prod_{a=1}^r\mathbb C\{u_a\}\ \times\mathcal O^{\varepsilon},
 \quad
 \iota:A\longrightarrow B,\quad
 \xi\longmapsto
       (E_a(t)+u_aR_a(t))_{a=1}^r\ \times\theta_0(t).
 \tag{HC5}
\]
The last component is omitted if $\varepsilon=0$. This is a
unital $\mathcal O$-algebra map in the original coordinate.
Each $Q_a$ is irreducible over $\mathcal K$: its discriminant
$4tR_a^2$ has odd $t$ valuation and cannot be a square in
$\mathcal K$. The factors are distinct because their leading
root slopes have different squares. The linear factor, if present,
is different from every quadratic factor. Therefore evaluation
into the product of the fraction fields is injective modulo $W$,
which proves injectivity of $\iota$.

The component map is onto. Indeed
$u_a=(\xi-E_a(t))/R_a(t)$ in the quotient by $Q_a$, and
every convergent series in $u_a$ is uniquely $f(t)+u_ag(t)$.
The central component is $\mathcal O$ itself. Consequently $B$
is finite free of rank $m$ over $\mathcal O$, and is also finite
as an $A$-module. It is the integral closure of $A$ in its product
of fraction fields. To verify that assertion directly, each factor
$\mathbb C\{u_a\}$ is a discrete valuation ring: every nonzero
series is a power of $u_a$ times a unit. A Laurent series of negative
valuation cannot satisfy a monic equation with coefficients of
nonnegative valuation, because its highest power would be the unique
term of smallest valuation. This proves integral closedness, also
for the central factor. An element integral over $A$ is integral
over $\mathcal O$ by finiteness and therefore has nonnegative
valuation in every factor. It belongs to $B$. Conversely finiteness
of $B$ over $A$ proves integrality by the determinant argument for
multiplication on a finite generating module. These two inclusions
prove the assertion.

The original algebra $A$ is retained throughout; $B$ is an auxiliary
finite algebra with the displayed inclusion and a computed quotient.

\HC\textbf{Every elementary divisor of the original inclusion.}
Use the basis of $B$ consisting first of its component identities
(including the central one), then $u_1,\ldots,u_r$. Let $J(t)$
be the matrix of HC5 from the ordered original basis
$1,\xi,\ldots,\xi^{m-1}$. Its entries are explicitly
\[
 \begin{split}
 J_{a,k}^{\rm even}(t)
  &=\frac{\theta_a(u)^k+\theta_a(-u)^k}{2},\\
 J_{a,k}^{\rm odd}(t)
  &=\frac{\theta_a(u)^k-\theta_a(-u)^k}{2u},\\
 J_{0,k}(t)&=\theta_0(t)^k\quad(\varepsilon=1),
       \qquad 0\le k<m.
 \end{split}
 \tag{HC6}
\]
These formulas retain the full original branches. Their precise
divisibility is sufficient to give the factorization
\[
 J(t)=C(t)D(t),\qquad
 D(t)=\operatorname{diag}_{0\le k<m}t^{\lfloor k/2\rfloor},
 \qquad C(t)\in\operatorname{GL}_m(\mathcal O).
 \tag{HC7}
\]
Here $C$ is defined entry by entry by HC6 divided by the stated
column power of $t$; it is not omitted from any metric or map.

For proof, $\theta_a(u)$ has order one, so the even part of its
$k$th power is divisible by $u^{2\lceil k/2\rceil}$ and its odd
part divided by $u$ by $u^{2\lfloor k/2\rfloor}$. The central
branch has order at least one in $t$. Thus all entries defining
$C$ are analytic. To check its determinant, order the even columns
before the odd columns for this check only. For $k=2\ell$ the even
rows at $t=0$ are $b_a^{2\ell}$, and the central row is $1$ for
$\ell=0$ and $0$ otherwise. The odd rows in these columns are
$4\ell c\,b_a^{2\ell-1}$, with value zero when $\ell=0$.
For $k=2\ell+1$ the even and central rows vanish and the odd
rows are $b_a^{2\ell+1}$. The diagonal blocks are consequently
Vandermonde matrices on the distinct numbers $b_a^2$, together
with $0$ in the even block when the central branch is present;
the odd block has the extra nonzero row factors $b_a$.
Both blocks are square and invertible. Thus $\det C(0)\ne0$,
proving HC7 in the original column order as well.

Multiplication by $C^{-1}$ induces the exact module isomorphism
\[
 \boxed{B/A\simeq
   \bigoplus_{k=0}^{m-1}\mathcal O/(t^{\lfloor k/2\rfloor}),
 \qquad
 \operatorname{length}_{\mathcal O}(B/A)
       =\left\lfloor\frac{(m-1)^2}{4}\right\rfloor.}
 \tag{HC8}
\]
The summand with exponent zero is the zero module. The length is
$r(r-1)$ for $m=2r$ and $r^2$ for $m=2r+1$, by summing the
exponents. HC7, rather than a rank comparison alone, proves every
elementary divisor.

\HC\textbf{The conductor with its full arithmetic factors.}
The conductor is the ideal
$\mathfrak c=\{b\in B:bB\subseteq A\}$, regarded in $B$ via HC5.
For any component $a$ let $W_a=W/Q_a$ and evaluate at its original
branch; for the central component use $W_0=W/Q_0$. Then
\[
 \begin{split}
 \mathfrak c_a&=W_a(t,\theta_a(u_a))\mathbb C\{u_a\}
            =u_a^{\,2(r-1)+\varepsilon}\mathbb C\{u_a\},\\
 \mathfrak c_0&=W_0(t,\theta_0(t))\mathcal O
            =t^r\mathcal O\quad(\varepsilon=1).
 \end{split}
 \tag{HC9}
\]
The first equalities retain the exact unit factors as well as the
orders. The conductor is the product of these component ideals.

To prove it, an element supported in just component $a$ belongs to
$A$ precisely when it is the image of a polynomial divisible by
all the other factors $Q_b$. The ring $\mathcal O[\xi]$ is a
unique factorization domain (polynomial Gauss lemma over the discrete
valuation ring); its distinct monic irreducible factors therefore
have product $W_a$ dividing that polynomial. Conversely every
multiple of $W_a$ vanishes in the other components. The projection
$A\to B_a$ is onto by HC5, so the ideal of possible $a$-components
is exactly $W_a(\theta_a)B_a$. Multiplying by the component
idempotents of $B$ proves both necessity and sufficiency for the
conductor definition.

For $b\ne a$, HC1 gives
\[
 Q_b(t,\theta_a(u_a))
      =(b_a^2-b_b^2)u_a^2+O(u_a^3),\qquad
 Q_0(t,\theta_a(u_a))=b_a u_a+O(u_a^2).
 \tag{HC10}
\]
Every displayed leading coefficient is nonzero. On the central
component $Q_b(t,\theta_0(t))=-b_b^2t+O(t^2)$. Multiplying gives
exactly the orders in HC9. For $m=1$ this gives conductor $\mathcal O$;
for $m=2$ it gives $B=A$. All components and all multiplicities
are included.

\HC\textbf{The special-fibre map and the loss measured by its kernel.}
At $t=0$ the inclusion has the explicit specialization
\[
 \bar\iota:\mathbb C[\xi]/(\xi^m)
   \longrightarrow
      \prod_{a=1}^r\mathbb C[u_a]/(u_a^2)\times\mathbb C^\varepsilon,
 \qquad \xi\longmapsto(b_au_a)_a\times0.
 \tag{HC11}
\]
For $m\ge2$ its kernel is exactly $(\xi^2)$ and its image has
dimension two, spanned by the images of $1,\xi$. Its cokernel
has dimension $m-2$. For $m=1$ it is an isomorphism.
Indeed powers $\xi^k$ with $k\ge2$ map to zero, while the two
displayed vectors are independent because some $b_a\ne0$.
This proves the assertions directly and records the precise failure
of injectivity after special-fibre base change. No nilpotent is
silently removed from the original $A/tA$.

The trace form also pulls back exactly. For real $t$ the coefficient
involution fixes $\xi$ and conjugates coefficients. On $B$ in the
stated component basis its Hermitian trace matrix is
\[
 G_B(t)=\operatorname{diag}
      (\underbrace{2,\ldots,2}_{r},
        \underbrace{1}_{\varepsilon},
        \underbrace{2t,\ldots,2t}_{r}),
 \qquad G_A(t)=J(t)^*G_B(t)J(t).
 \tag{HC12}
\]
The entry $1$ is omitted when $\varepsilon=0$. In a quadratic
component multiplication by $u_a$ has matrix
$\left(\begin{smallmatrix}0&t\\1&0\end{smallmatrix}\right)$;
its trace is zero and the trace of $u_a^2$ is $2t$. This proves
$G_B$. Over $\mathcal K$ the inclusion is an algebra isomorphism,
so regular traces agree there; both sides of HC12 are analytic,
giving the equality also at zero. In particular
\[
 G_A(0)=\operatorname{diag}(m,0,\ldots,0).
 \tag{HC13}
\]
Thus the entire nilpotent radical of the original local Weil trace
EP11 is retained. The appearance of several dual-number factors
in the auxiliary special fibre does not replace that original trace.

\HC\textbf{The discriminant, including its first arithmetic correction.}
Let $\mathfrak D_A(t)$ denote the determinant of the regular
bilinear trace Gram matrix in $1,\xi,\ldots,\xi^{m-1}$.
At real $t$ it is also the determinant in HC12. Put
$M=m(m-1)/2$ and $b=2d-c^2$. Then
\[
 \boxed{\mathfrak D_A(t)=D_m t^M
      \bigl(1+m(m-1)b\,t+O(t^2)\bigr),\qquad
 D_m=2^M\prod_{j=1}^m j^j.}
 \tag{HC14}
\]
The empty-product conventions give $D_1=1$ and a constant
discriminant for the one-dimensional algebra.

To prove the full coefficient, evaluate the original power basis at
all the $m$ branches. Its Vandermonde determinant squared is
\[
 \mathfrak D_A(t)=\prod_{i<j}
                    (\theta_i(u)-\theta_j(u))^2.
 \tag{HC15}
\]
HC1 makes every gap $u(r_i-r_j)(1+bu^2)+O(u^4)$.
The common $2cu^2$ translation cancels. Consequently the product
is $u^{2M}\prod_{i<j}(r_i-r_j)^2(1+2Mb u^2+O(u^3))$.
It is invariant under $u\mapsto-u$ because that permutation
preserves the entire multiset of branches. Its analytic expansion
in $t$ therefore improves the remainder to $O(t^2)$ after removing
$t^M$.

For the constant, the monic Hermite recurrence proved in EP5 is
$P_n'=nP_{n-1}$ and $P_n=zP_{n-1}-2(n-1)P_{n-2}$.
Define $R_n=\operatorname{Res}(P_n,P_{n-1})$ by products over roots.
Exchanging the two monic factors introduces sign
$(-1)^{n(n-1)}=1$. Evaluation at the roots of $P_{n-1}$ gives
$R_n=(-2(n-1))^{n-1}R_{n-1}$, with $R_1=1$.
The derivative product formula for a discriminant then gives
$\operatorname{disc}P_n=(-1)^{n(n-1)/2}n^nR_n
=2^{n(n-1)/2}\prod_{j=1}^n j^j$. This proves HC14.
The Hermite convention and human-source attribution remain precisely
\href{https://dlmf.nist.gov/18.5.E13}{DLMF18.5.13}, with its authors
and explicit coordinate conversion given in EP3.

HC12 also gives
$\mathfrak D_A=(\det J)^2\,2^{2r}t^r$.
Taking orders and using HC8 proves the exact identity
\[
 M=r+2\operatorname{length}_{\mathcal O}(B/A).
 \tag{HC16}
\]
Thus the discriminant order separates into the branch ramification
order and twice the explicitly computed conductor quotient length.

The same coefficient is measured by the original cluster energy:
\[
 \left.\frac{d}{dt}\log
       \frac{\mathfrak D_A(t)}{D_m t^M}\right|_{t=0}
   =m(m-1)(2d-c^2)
   =-4\operatorname{FP}_{t=0}E_{x_0}(t)
   =-\frac{m(m-1)}4(\log u_\rho)''(0).
 \tag{HC17}
\]
The quotient inside the logarithm is an analytic unit with value
one, so this derivative requires no choice of logarithm.
The two last equalities follow by substituting EP8 and EP15--16;
they retain the original arithmetic local unit, not a substituted
polynomial source.

\HC\textbf{The support lift retains every connecting map.}
For any of the original bounded distributive support lattices $L$,
apply $G_L$ to HC5. The map
\[
 G_L(A)\longrightarrow G_L(B),\quad
 (a,1)\longmapsto(\iota(a),1),\qquad
 (0,\lambda)\longmapsto(0,\lambda)
 \tag{HC18}
\]
is injective: its amplitude map is injective and it fixes every
support label. Addition and multiplication are preserved by the
original definitions. Its arithmetic reflection is exactly $\iota$.
Under the unital localization at supported $e$ both sides map to
$L$, and the induced map on $L$ is the identity. This follows
directly from the universal localization proved in SH: every
top-supported arithmetic value becomes the top unit, while every
zero-support label remains. The full diagram commutes:
\[
 \begin{tikzcd}[column sep=large]
 G_L(A)\arrow[r,hook]\arrow[d,"p_A"']
  &G_L(B)\arrow[d,"p_B"]\\
 A\arrow[r,hook,"\iota"']&B
 \end{tikzcd}
 \tag{HC19}
\]
The support receivers commute in the separate diagram:
\[
 \begin{tikzcd}
 G_L(A)\arrow[r,hook]\arrow[d,"\chi_A"']&
 G_L(B)\arrow[d,"\chi_B"]\\
 L\arrow[r,equal]&L.
 \end{tikzcd}
\]
The quotient HC8, conductor HC9 and trace HC12 are computed before
this amplitude collapse. The support localization cannot reconstruct
them from its value alone; their exact receiver is the retained
arithmetic reflection, which does not extend through that localization.
For example $e$ acts as zero on the arithmetic $A$-module $B/A$,
so its base change along the localization is zero: on a tensor
$1\otimes v$, the equality $1=e$ gives
$1\otimes v=1\otimes ev=0$.

\HC\textbf{The relative cotangent defect and its exact endpoint.}
The full time family has the two-term cotangent complex
\[
 L_{A/\mathcal O}=[A\xrightarrow{\ W_\xi\ }A\,d\xi],
       \qquad\text{degrees }-1,0.
 \tag{HC20}
\]
One construction starts with the free graded algebra over
$\mathcal O[\xi]$ on one exterior generator $\eta$ of degree $-1$,
with differential $d\eta=W$. Its underlying two-term complex is
exact in degree $-1$ because multiplication by the monic $W$ is
injective; its degree-zero cokernel is $A$. Taking relative
differentials and tensoring with $A$ sends $d\eta$ to
$W_\xi d\xi$, proving HC20.

Multiplication by $W_\xi$ is injective on $A$: in the product of
fraction fields of HC5 its value on each branch is the product
of all nonzero differences from the other branches. It is invertible
there, and $A$ injects into that product. Thus
\[
 H^{-1}(L_{A/\mathcal O})=0,\qquad
 \Omega=H^0(L_{A/\mathcal O})=A/(W_\xi)\,d\xi.
 \tag{HC21}
\]
The module $\Omega$ is a finite torsion $\mathcal O$-module with
\[
 \operatorname{length}_{\mathcal O}\Omega=M=\frac{m(m-1)}2.
 \tag{HC22}
\]
For proof, the determinant of multiplication by $W_\xi$ in $A$
is the product of its branch values, hence
$(-1)^M\mathfrak D_A(t)$. A square matrix over the discrete
valuation ring with nonzero determinant has cokernel length equal
to the determinant valuation: successive row and column elimination
selects an entry of least valuation, clears its row and column by
division, and iterates to a diagonal matrix of powers of $t$ times
units. Its cokernel has the corresponding sum of exponents, equal
to the valuation of the determinant. HC14 now proves HC22.

Specializing the actual complex, rather than its cohomology alone,
gives $[A_0\xrightarrow{m\xi^{m-1}}A_0d\xi]$, where
$A_0=\mathbb C[\xi]/\xi^m$. The free $\mathcal O$-resolution
in HC21 gives the exact base-change calculation
\[
 \begin{split}
 \operatorname{Tor}_1^{\mathcal O}(\Omega,\mathbb C)
    &\simeq(\xi)\subset A_0,\\
 \Omega\otimes_{\mathcal O}\mathbb C
    &\simeq A_0/(\xi^{m-1})\,d\xi.
 \end{split}
 \tag{HC23}
\]
For $m=1$ both sides vanish. For $m\ge2$ both displayed spaces
have dimension $m-1$. The first map sends the cycle represented
by $a\eta$ to its original coefficient $a\in(\xi)$.
Consequently the endpoint radical of the original Weil trace
EP11 is precisely this explicit derived base-change group.
It occurs even though the relative family has zero $H^{-1}$.
This extends the programme's collision Jacobian/cotangent
comparison \href{https://github.com/KokunoYumeto/zeta-function-research-reader/blob/5a89872598df0902b7c1393cf8e4692ca3010a95/workbenches/splitzero-tandem/branches/tau-zero-prime-positivity/COLLISION_RESIDUE_DUALITY_DERIVATION.md}{RD27--28a}
to the original arithmetic time family;
the complete argument is given here, with its original coordinates.

\HC\textbf{A perfect observation of the torsion throughout the family.}
For $f\in A$ let $\lambda_A(f)$ be the coefficient of
$\xi^{m-1}$ in its unique representative of degree less than $m$.
The bilinear form $\lambda_A(fg)$ is perfect over $\mathcal O$.
In the power basis its entries with $i+j<m-1$ are $0$
and those with $i+j=m-1$ are $1$: products of degree less than
$m$ require no reduction. Reversing the columns therefore gives
a triangular matrix with diagonal $1$. Its determinant is
$(-1)^M$, a unit. This proves perfection without deleting the
terms above the anti-diagonal.

The exact trace identity is
\[
 \operatorname{Tr}_{A/\mathcal O}(M_f)
       =\lambda_A(W_\xi f).
 \tag{HC24}
\]
On a nonzero fibre, Lagrange interpolation says
$\lambda_A(g)=\sum_j g(\theta_j)/W_\xi(\theta_j)$:
the $j$th interpolation basis polynomial has leading coefficient
$1/W_\xi(\theta_j)$. Taking $g=W_\xi f$ gives HC24 there.
Both sides are analytic in $t$, so the identity holds at zero too.
This proves the trace-to-dual factorization through the actual
Jacobian multiplier $W_\xi$.

Omit the symbol $d\xi$ when identifying $\Omega$ with its coefficient
module in the following explicitly defined pairing:
\[
 \mathscr L:\Omega\times\Omega\longrightarrow\mathcal K/\mathcal O,
 \qquad([a],[b])\longmapsto
      \lambda_A\left(\frac{ab}{W_\xi}\right)\pmod{\mathcal O}.
 \tag{HC25}
\]
The inverse of $W_\xi$ is taken in $A\otimes\mathcal K$, where
it exists by HC21. Adding $W_\xi c$ to either representative
changes the value by $\lambda_A(cb)\in\mathcal O$, proving
well-definedness. If the value is zero for all $b\in A$, perfection
of $\lambda_A$ says that $a/W_\xi$ belongs to $A$, and therefore
$a\in W_\xi A$. This proves nondegeneracy. It is perfect: a
cyclic module $\mathcal O/(t^q)$ has dual into
$\mathcal K/\mathcal O$ equal to $t^{-q}\mathcal O/\mathcal O$,
of the same length $q$. Applying the elementary divisor argument
of HC22 to $\Omega$ shows that the injective duality map just
proved has equal-length source and target, hence is an isomorphism.

At the endpoint the corresponding finite duality is
\[
 (\xi)\times A_0/(\xi^{m-1})\longrightarrow\mathbb C,
           \quad(a,[b])\longmapsto\lambda_{A_0}(ab).
 \tag{HC26}
\]
It is well defined because $\xi\xi^{m-1}=0$. The bases
$\xi,\ldots,\xi^{m-1}$ and $1,\ldots,\xi^{m-2}$ have a
reversed identity pairing matrix, so it is perfect. Thus the
trace-invisible endpoint directions have both an exact family
torsion observation HC25 and an exact endpoint dual receiver HC26.
These are bilinear residue pairings, with the trace relation HC24;
they are not assigned positivity that their definitions do not give.

\HC\textbf{The full prime spectrum at the actual collision.}
The preceding calculation also determines the new prime fibres;
it must not be read only at the level of complex vector spaces.
Write $z_\lambda=(0,\lambda)$ and let
$\mathfrak e=(e)\subset G_L(R)$ for any ring $R$.
The multiplication $ez_\lambda=z_\lambda$ gives
\[
 \mathfrak e=\{z_\lambda:\lambda\in L\},\qquad
 \sqrt{\mathfrak e}=p^{-1}(\operatorname{Nil}R).
 \tag{HC27}
\]
Indeed every element of the principal ideal has zero amplitude,
every zero label is its displayed multiple, and a top-supported
$a$ has a power in this ideal exactly when $a$ is nilpotent.
For the original $A_0=\mathbb C[\xi]/(\xi^m)$ this is
\[
 \sqrt{(e)}=Q_{(\xi)}=p^{-1}((\xi)),
 \qquad V(e)=\{Q_{(\xi)}\}.
 \tag{HC28}
\]
When $m\ge2$, $(e)$ itself is not prime in this collision algebra:
the nonzero top-supported elements $\widehat\xi$ and
$\widehat{\xi^{m-1}}$ multiply to $e$, and neither belongs to $(e)$.
For $m=1$ the amplitude ring is a field and $(e)$ is prime.
This states precisely how the supported-zero prime in the integral
coefficient ring changes under the actual nonreduced specialization.
It retains the supported ideal, its radical and its prime receiver.

Here is the complete classification, including every lower label.
A semiring prime omitting $e$ contains no top-supported amplitude:
if it contained $\widehat a$, multiplying by $e$ would put $e$
in that prime. Its zero labels form a proper lattice prime ideal
$J$, and the corresponding semiring prime is
$P_J=\{z_\lambda:\lambda\in J\}$. Closure under addition is closure
under joins, and multiplicative primality is primality for meets.
Conversely these conditions verify directly that every $P_J$ is a
semiring prime. A prime containing $e$ contains every zero label;
its top amplitudes form a ring prime $\mathfrak p$, since addition,
multiplication and additive inversion (multiplication by $-1$)
are preserved. Its exact form is $Q_{\mathfrak p}=p^{-1}(\mathfrak p)$.
Conversely each such inverse image is prime. Every $P_J$ is contained
in every $Q_{\mathfrak p}$, because its elements have zero amplitude.
This reproves the complete support/arithmetic specialization order.

For the original special-fibre map HC11 the resulting spectral map is
\[
 \begin{array}{ccc}
 \operatorname{Spec}_{\rm lat}L\ \blacktriangleright\
       \{q_1,\ldots,q_r,q_0\text{ if }\varepsilon=1\}
 &\longrightarrow&
 \operatorname{Spec}_{\rm lat}L\ \blacktriangleright\{q\},\\
 P_J&\longmapsto&P_J,\\
 q_a\ \text{and }q_0&\longmapsto&q.
 \end{array}
 \tag{HC29}
\]
Here $q=Q_{(\xi)}$, and $q_a$ is the prime of the $a$th dual-number
factor (with all other product components in the ideal); $q_0$
is the central field-factor prime. Evaluation of HC11 proves every
inverse image in HC29. The coefficient object on the left is
$G_L(B_0)$ with its one complete support lattice $L$, not the
larger product of independent support carriers. Thus the exact
number $r+\varepsilon$ of arithmetic branch points, their common
image, the shared support primes and the conductor are all retained.
Localizing $e$ keeps the identity on the displayed lattice spectrum;
it removes the arithmetic points through the same universal map
as HC18--19. The branch and holonomy information is carried by
the original inclusion and its arithmetic fibres before that map.

\HC\textbf{The discriminant is measured by the holonomy commutator.}
Use the original coefficient holonomy $M_A(t)$ and the complete
scaling multiplier $S_L(t)$ of HH9 and HH29. For fixed real $L$
and small positive $t$, the reflected trace is positive definite,
so it defines the Hilbert--Schmidt norm on its endomorphisms.
The exact expansion proved in HH33--38, combined with HC17, gives
\[
 \begin{split}
 \|[M_A(t),S_L(t)]\|_{\rm HS,T_t}^2
  ={}&2L^2m(m-1)t\\
   &+\left[
       4L^2\left.\partial_t\log
           \frac{\mathfrak D_A(t)}{D_m t^M}\right|_{t=0}
       -\frac{L^4}3m(m-1)(2m-3)
     \right]t^2+O(t^3).
 \end{split}
 \tag{HC30}
\]
No metric is replaced: HH18 identifies this exact norm with the
Euclidean norm of value vectors at the same positive time, and
HC12 gives its full original-coordinate matrix. HC17 evaluates
the derivative by the original $a_n$. Thus the conductor
discriminant, the arithmetic energy finite part and the measured
noncommutativity of the two transports have an explicit common
coefficient. This is a quantitative receiver for the original
geometry; no sign of the complete global Weil form is inferred
from this positive-time local norm.

\begin{figure}[p]
\centering
\begin{tikzpicture}[>=Stealth,node distance=1.6cm,
 box/.style={draw,rounded corners,align=center,text width=11.5cm,inner sep=9pt}]
\node[box] (a) {Original collision algebra
 $A=\mathbb C\{t\}[\xi]/W$\\
 Full arithmetic branches, original $\xi=Z-x_0$, all coefficients retained};
\node[box,below=of a] (b) {Auxiliary integral closure
 $B=\prod_{a=1}^r\mathbb C\{u_a\}\times\mathbb C\{t\}^{\varepsilon}$\\
 $t=u_a^2$, $\xi=E_a(t)+u_aR_a(t)$};
\node[box,below=of b] (q) {Exact quotient:
 $B/A\simeq\displaystyle\bigoplus_{k=0}^{m-1}
                \mathcal O/(t^{\lfloor k/2\rfloor})$\\
 Length $\lfloor(m-1)^2/4\rfloor$;
 original special fibre $\mathbb C[\xi]/(\xi^m)$ retained};
\draw[->] (a)--node[right] {$\iota$, matrix $J=CD$}(b);
\draw[->] (b)--node[right] {quotient by the original image}(q);
\end{tikzpicture}
\caption{The full actual collision, not a replacement polynomial.
HC1--7 define every arrow using the original $H_t$. HC8 evaluates
the complete finite-module defect. HC9 retains every conductor
unit, HC11 gives the possibly noninjective special-fibre map, and
HC12--17 connect the discriminant to the original energy and
completed-zeta unit. The source heat equation is Rodgers--Tao v5;
the exact Hermite leading polynomial and its convention are cited
and proved in EP2--6.}
\end{figure}
